% !TeX root = ../教材题库_全部条目_转写工作稿.tex
% Source: XJ-B1-C04-EX-P132-Y1265
% 第4章 / 4.4 函数与方程 / 教材 p.132、133、134 / PDF 第139、140、141页
% OCR未校对：公式与文字均可能有误；未完成数学TeX转写；禁止直接发给学生
% 原题及相邻原解截图（核对用）：../题目原图/XJ-B1-C04-EX-P132-Y1265-01.webp, ../题目原图/XJ-B1-C04-EX-P132-Y1265-02.webp, ../题目原图/XJ-B1-C04-EX-P132-Y1265-03.webp
% 本文件只在转写完成后作为题干使用；OCR 条目尚不是可靠题干。
\begin{bankitem}{4.4 例题 例题4}{第4章 / 4.4 函数与方程 / 教材 p.132、133、134 / PDF 第139、140、141页}{XJ-B1-C04-EX-P132-Y1265}{OCR未校对：公式与文字均可能有误；未完成数学TeX转写；禁止直接发给学生}
\begin{ocrdraft}
[教材p.132 / PDF第139页]
\par
0二4]在例3中，我们已经说明了不到一一3十1在[0，0上恰有一个零
\par
点．试用二分法来计算这个零点的更精确的近似值（误差不超过0．01).
\par
解前面已经知道函数丆0）的零点在[0，0上，此时包含零点的区间长度为
\par
1一0一1，远超0．01，不符合要求．于是我们可以通过二分法来缩小零点所在区间，
\par
具体如下：
\par
0十1
\par
第1次，取[0，0的中点
\par
=0．5，用计算器或计算机求出不0．5）0，38>0．
\par
2
\par
由不1)一一1<0，可知不0．5）．不D<0，于是零点在[0．5，0上，此时包含
\par
零点的区间长度为1一0．5=0．5，远超0，01，不符合要求．
\par
[教材p.133 / PDF第140页]
\par
0．5十1
\par
第2次，取[0．5，0的中点
\par
=0.75，求出不0．75）
\par
—0.27<0．
\par
2
\par
由于不0．5）．不0，75）<0，可知零点在[0．5，0，75]上，此时包含零点的区间
\par
长度为0．75一0，5=0．25>0．01，不符合要求．
\par
继续重复上述步骤，可得下表：
\par
不m）的
\par
零点所在
\par
是否符合
\par
a+b
\par
次数
\par
区间左端点a区间右端点方
\par
2
\par
近似值
\par
区间长度
\par
要求
\par
1
\par
0
\par
1
\par
0．38
\par
否
\par
2
\par
1
\par
0．75
\par
一0．27
\par
0．25
\par
否
\par
3
\par
0．75
\par
0．625
\par
0．07
\par
0．125
\par
否
\par
4
\par
0．625
\par
0．75
\par
0．6875
\par
一0．09
\par
0．0625
\par
否
\par
否
\par
5
\par
0．625
\par
0．6875
\par
0．65625
\par
一0．009
\par
0．03125
\par
6
\par
0．625
\par
0．65625
\par
0．640625
\par
0．032
\par
0．015625
\par
否
\par
7
\par
0．640625
\par
0．65625
\par
0．6484375
\par
0．011
\par
0．0078125
\par
是
\par
由上表可知，计算到第7次，由于不0．6484375）\ensuremath{\cdot}不0．65625）<0，可知包
\par
含零点的区间为匚0．6484375，0．65625〕，零点所在区间长度小于0，01.于是，该
\par
区间内任意一点都可作为零点的近似值，例如可取0，65（也可取区间的端点值）作为
\par
零点的近似值，也即可取0．65作为方程一32十1一0的一个近似解．
\par
从上述计算过程，可以看出用二分法求函数零点近似值的一般操作方法．
\par
设函数3，一不x）定义在区间D上，其图象是一条连续曲线．我们希望求它在D
\par
上的一个零点“的近似值使它与零点的误差不超过给定的正数\ensuremath{\text{[OCR:EUR]}},即使得
\par
第1步，在D内取一个闭区间，国匚D，使丆（a）与不b）异号，即f(a)
\par
不<0．
\par
第2步，取区间@，刿的中点
\par
m=—(a十b).
\par
第3步，计算f(m)，并进一步确定零点所在的区间：
\par
（1）若不m）一0（此时一”D，则m就是函数的零点；
\par
（2）若不磊）．不m）<0（此时e，m)），则令b=m；
\par
（3）若f(m)\ensuremath{\cdot}不<0（此时e(m，），则令磊一．
\par
第4步，判断是否达到精确度：若汤一，则得到零点的近似值或仞，
\par
也可选取区间内任意一个数作为零点的近似值；否则重复步骤2、4．
\par
由于用二分法求函数零点近似值时，通常计算量大且需重复相同的步骤，因而
\par
[教材p.134 / PDF第141页]
\par
可以借助计算机通过设计如下程序（图4．4\_5）进行解决．
\par
开始
\par
定义丿、（到
\par
输入8，b
\par
a+b
\par
2
\par
否
\par
丿、间\ensuremath{\cdot}厶<0
\par
否
\par
否
\par
感兴趣的司学，可以
\par
在此基詘上利用信息技术
\par
求解例4．
\par
输出解无：口
\par
结束
\par
图4．4一5
\end{ocrdraft}
\end{bankitem}
